% ============================================================
%  CHAPTER 1 — INTRODUCTION
%  Sections: 1.1 Introduction | 1.2 Problem Statement | 1.3 Academic Context
% ============================================================

\chapter{Introduction}
\label{chap:introduction}

% ============================================================
\section{The Relevance of Corner Kicks in Modern Football}
\label{sec:intro_relevance}
% ============================================================

Football is a sport of marginal gains. Across the top five
European leagues, the average match is decided by a single goal, and any
repeatable, high-frequency situation that shifts the probability of scoring, even marginally, carries strategic value that far exceeds its apparent
simplicity. Corner kicks are among the clearest expressions of this logic. They
occur, on average, 5.5 times per team per match \citep{Casal2015}, they restart
play in the most dangerous corridor on the pitch, and --- unlike almost any
phase of open play --- they permit prior rehearsal of exact player positions and
movement sequences. A professional side will typically execute more than two
hundred corners across a thirty-four match season: two hundred scripted
opportunities to attack, from a fixed point, in front of the opponent's goal.

Yet the raw conversion rate is strikingly low. Only about 2--3\,\% of corners
across elite competitions result directly in a goal \citep{Casal2015,
ShawGopaladesikan2020}. This tension between high frequency and low conversion
has long invited the question of whether corners are genuinely worth investing
in, or whether they are, as Power et al.\ famously tested, a tactical dead end
\citep{Power2018}. The answer that data have consistently returned is the former.
The per-possession probability of scoring from a corner is approximately 1.8\,\%,
compared with only 1.1\,\% for an open-play possession --- a 64\,\% relative
increase \citep{Power2018}. More strikingly, when a goal is scored from a corner,
it is decisive --- winning or drawing the match --- in roughly three quarters of
cases \citep{Casal2015}. And the aggregate weight of corner goals in the overall
goal distribution has grown substantially over the past decade: in the English
Premier League during the 2023/24 season, corners directly contributed to
approximately 16\,\% of all non-penalty goals, continuing a trend that has
accelerated in recent years as clubs have invested more deliberately in set-piece
preparation.

The realisation that corners are a systematically underexploited source of goals
has driven a structural change in professional club football. Dedicated
\emph{set-piece coaches} --- specialists whose sole responsibility is the design,
preparation, and analysis of dead-ball situations --- have become near-universal
at the elite level. The career trajectory of Nicolas Jover, who began as a video
analyst at Montpellier, became set-piece coach at Brentford, then Manchester
City, and eventually Arsenal, illustrates the institutionalisation of the role.
Under Jover's guidance at Arsenal, the club scored twenty set-piece goals in the
2023/24 Premier League season, the most in the division, and broke or equalled
multiple Premier League records for corner effectiveness. Italian specialist
Gianni Vio, who introduced data-driven corner routines at Brentford as early as
2015/16, saw the club's corner-to-goal conversion rate improve by 80\,\% in a
single season. FC Midtjylland, another data-forward club operating under Matthew
Benham's analytical approach, scored 25 set-piece goals in one Danish Superliga
campaign --- fourteen more than the next best side.

These are not isolated success stories. They reflect a general recognition that
set pieces, and corners in particular, constitute a \emph{market inefficiency}:
an area of the game where the expected return on investment in time, analysis,
and tactical preparation substantially exceeds the return on equivalent
investment in open-play development \citep{Power2018}. Paul Power, co-author of
the 2018 Sloan set-piece paper, has argued publicly that a high-quality set-piece
programme is, in goal-expectation terms, roughly equivalent to spending
\pounds80 million on a world-class striker. Corner kicks, in short, matter ---
and the question of how to exploit them more effectively, and how to prepare
against specific opponents doing the same, is one of the most practically
consequential open problems in applied football analytics.

% ============================================================
\chapter{Problem Statement}
\label{sec:problem_statement}
% ============================================================

\section{The Preparation Gap}
\label{sec:prob_preparation}

Despite the growing recognition of corners as a strategic asset and the
proliferation of specialist coaching roles, the \emph{preparation process} for
corners has not evolved at the same pace as the ambition surrounding them.
The dominant workflow in professional clubs today remains notational video
analysis: an analyst watches the opposition's previous matches, classifies each
of their corners by hand --- the delivery type, where the ball landed, how the
defending team organised itself, which attackers ran where --- and synthesises
these observations into a pre-match scouting report. Wright et al.\ (2013)
estimated that performance analysts spend, on average, approximately six hours
per match preparing opposition analysis; set-piece review accounts for a
substantial portion of this time. Yet Gianni Vio, speaking from two decades of
specialised set-piece practice, has observed that despite corners and free kicks
accounting for approximately 30\,\% of all goals at elite level, clubs typically
dedicate no more than ten minutes per week of training time to set-piece work.
The gap between the \emph{strategic importance} of the set piece and the
\emph{analytical bandwidth} available to exploit that importance is, by
practitioners' own accounts, enormous.

This gap has a structural cause that goes beyond the number of hours available.
Manual notational analysis of corners is not only slow; it is also
fundamentally limited in what it can observe. A human analyst watching a video
frame can categorise whether a defender \emph{appears} to be following an
attacker or holding a zone, but cannot precisely quantify the distance between
those two players at each moment leading up to delivery, the velocity with which
the defender is tracking, the degree of synchronisation between a defender's
movement and an attacker's run, or the spatial coverage gaps that result from
the collective configuration of all ten outfield defenders. These are
\emph{continuous, spatiotemporal} properties of the scene, and they cannot be
adequately captured in the discrete, categorical labels that notational analysis
produces. The consequence is that much of the tactical information embedded in a
corner --- information that is, in principle, available in the data --- is lost
in the process of manual observation.

\section{The Scalability Problem}
\label{sec:prob_scalability}

The problem is compounded by scale. A team preparing for a specific opponent
will typically have access to no more than ten to fifteen of that opponent's
most recent matches, yielding perhaps sixty to ninety corners to work with.
Because short corners, anomalous situations, and in-match tactical adjustments
reduce the homogeneous sample further, an analyst may have as few as forty
``standard'' defensive corners from a given opponent to draw conclusions from.
This is a statistically small sample from which to infer stable strategic
tendencies, particularly for the kind of conditional analysis that coaches need:
not simply ``how does this opponent defend corners?'' but ``how does this
opponent defend corners of the \emph{type we tend to take}, with our
\emph{specific personnel}?'' The manual pipeline, which is already slow for
single-team aggregate analysis, is wholly unsuited to the conditional,
opponent-specific, player-specific questions that coaches actually ask.

Tracking data offers, in principle, a path out of this bottleneck. Where a club
has access to a league-wide tracking database --- as is increasingly the case
through commercial data providers --- thousands of corners across a full season
of play can be processed simultaneously, yielding far larger samples per
analytical category, far more stable estimates of tendencies, and far richer
conditional analyses. The computation is also consistent: unlike a human analyst
who may code a given corner differently depending on the viewing angle, the
fatigue level of the session, or the ambiguity of a specific player's movement,
an algorithmic pipeline applied to the same tracking data will always produce
the same output. This consistency is particularly important when the goal is not
just to describe a single opponent's behaviour but to build a general model of
which attacking tactics work against which defensive configurations.

\section{The Two-Sided Blindspot}
\label{sec:prob_blindspot}

The academic and applied analytics literature has begun to engage with this
problem, but has done so in a piecemeal way. The most important prior work ---
Shaw and Gopaladesikan \citeyearpar{ShawGopaladesikan2020}, Bauer, Anzer, and
Smith \citeyearpar{Bauer2022}, and Wang et al.\ \citeyearpar{Wang2024} --- has
addressed either the defensive side (classifying how opponents organise
themselves to defend corners) or the offensive side (identifying what attacking
routines teams use and modelling who is likely to receive the ball) in
isolation. No published study has systematically combined a data-derived
characterisation of both the attacking routine and the defensive structure into
a single framework, and then asked the conditional question that motivates
set-piece coaching in the first place: \emph{which attacking configurations are
most effective against which defensive configurations?}

This two-sided blindspot has a direct practical consequence. A coach who knows
that a specific opponent defends corners predominantly with man-markers still
cannot, from that knowledge alone, determine which of his own attacking routines
is most likely to exploit that structure --- because the effectiveness of an
attacking routine against a man-marking defence depends on details (blocking
movement, run timing, delivery angle) that are themselves uncharacterised in
quantitative terms. Conversely, a coach who knows that inswinging deliveries
to the central penalty area tend to generate more xG on average cannot determine
whether that general finding applies when the opponent in question is defending
with a particular hybrid structure that systematically overloads that zone.
The missing link is an integrated, data-derived model that bridges the offensive
and defensive characterisations and operates at the level of specificity that
coaching decisions require.

\section{The Individual Quality Dimension}
\label{sec:prob_player}

A further limitation of the existing literature concerns the treatment of
individual players. Corner effectiveness is determined not only by which tactical
setup is deployed and against which defensive structure, but by \emph{who}
executes it. The quality of the corner taker --- measured by delivery precision,
trajectory consistency, and the xA generated by his deliveries --- determines
whether the ball reaches the intended zone at all. The aerial ability of the
target player --- measured by aerial duel win rate, height, timing, and opponent-
adjusted quality metrics --- determines whether a delivery in the right zone
produces a shot. These individual contributions are not separable from the
tactical analysis; a routine that performs well on average may perform very
differently when executed by a technically limited taker, or when the primary
aerial target is matched against a dominant defender. Yet the existing
quantitative literature on corners treats players as interchangeable within their
tactical roles, and does not model the interaction between player quality and
tactical setup. This is a gap that becomes increasingly important as the analysis
moves from the population level (what works on average?) to the decision level
(what should \emph{this} team do, with \emph{these} players, against
\emph{this} opponent?).

\section{Research Questions}
\label{sec:research_questions}

The preceding problem analysis motivates the central research question of this
thesis:

\begin{quote}
  \textbf{Main Research Question:} \emph{To what extent can tracking and event
  data be used to tailor attacking corner tactics to specific defensive
  structures, in order to maximise corner success?}
\end{quote}

This question is operationalised through five sequential sub-questions, each
of which corresponds to one analytical layer in the methodology:

\begin{description}
  \item[SQ1] \emph{To what extent can individual defensive roles be automatically
    classified from tracking data, and can these classifications be aggregated
    to characterise a team's overall defensive corner strategy?}

  \item[SQ2] \emph{Can attacking corner patterns be automatically classified from
    tracking and event data, and what taxonomy of offensive tactics emerges from
    this classification?}

  \item[SQ3] \emph{Which attacking corner strategies are most effective against
    which defensive structures?}

  \item[SQ4] \emph{Which individual player qualities and attributes influence the
    success of corners?}

  \item[SQ5] \emph{How can these combined insights be used to improve preparation
    for corners against specific opponents?}
\end{description}

SQ1 and SQ2 establish the two descriptive vocabularies --- one for defensive
setups, one for attacking routines --- that make SQ3 possible. SQ3 is the
central analytical contribution: the conditional effectiveness analysis that
links the two sides of the corner. SQ4 enriches the SQ3 findings by
disentangling tactical from individual effects. SQ5 synthesises the outputs
of SQ1--SQ4 into a practical preparation tool for coaching staff. The analytical
pipeline is described in full in Chapter~\ref{sec:methods}.

% ============================================================
\chapter{Academic Context}
\label{sec:academic_context}
% ============================================================

The research questions positioned above sit at the intersection of three
established bodies of academic and applied literature: (i) the notational
performance-analysis tradition on corner kicks, (ii) the coaching-science
literature on defensive and offensive set-piece systems, and (iii) the
spatiotemporal analytics literature that applies tracking data and machine
learning to football. A fourth, smaller body of work on individual player
attributes and aerial duel performance is also relevant, particularly for SQ4.
The review below works through each tradition in turn before identifying the
specific gap that this thesis addresses.

% ------------------------------------------------------------
\section{Notational Performance Analysis of Corner Kicks}
\label{sec:lit_notational}
% ------------------------------------------------------------

The largest body of academic work on corner kicks belongs to the notational
performance-analysis tradition. Studies in this tradition manually code corners
from video using structured observation instruments --- typically purpose-built
software platforms such as LINCE, Nacsport, or Wyscout Scout --- and examine
statistical relationships between coded categorical variables and outcomes such
as shots, shots on target, or goals. The observational units are human-assigned
labels: the side from which the corner is taken, the kicker's foot, the delivery
trajectory (inswing, outswing, or straight), the target zone in the penalty
area, the number of attackers and defenders in the box, the defensive system
employed, and the match context \citep{Sarmento2014, Sarmento2025}.

\subsubsection{Frequency, Conversion, and Match Impact}

Casal and colleagues \citeyearpar{Casal2015} provide the most frequently cited
quantitative baseline. Analysing 1{,}139 corners from 124 matches at the 2010
FIFA World Cup, UEFA Euro 2012, and the 2010/11 UEFA Champions League, they
reported a direct goal-per-corner rate of 2.2\,\%. Crucially, they found that
76\,\% of those goals were directly responsible for the scoring team winning or
drawing the match in which they were scored --- a figure that underlines the
match-decisive weight of corner conversions. Their multivariate analysis
identified three variables as significantly associated with scoring: time of
match ($p = 0.04$), number of intervening attackers ($p = 0.001$), and the
degree of offensive organisation ($p = 0.02$). Later work by
Fern\'andez-Herm\'ogenes, Camerino, and Hileno \citeyearpar{FernandezH2021},
working on 2{,}029 corners from the Spanish La Liga and Segunda Divisi\'on,
found that only 29.1\,\% of corners generated a shot attempt --- consistent
with the earlier range of 17--24\,\% shot-per-corner rates reported across
multiple studies and competitions \citep{Borras2005, SainzdeBaranda2012,
SanchezFlores2012}. Corner kicks occur roughly ten times per match at elite
level \citep{FernandezH2021}, providing a large base of observations for
statistical analysis while also highlighting the low per-event conversion
rate that makes each marginal improvement in targeting so valuable.

\subsubsection{Delivery Characteristics and Target Zones}

A parallel strand of the literature has focused on \emph{delivery
characteristics}. Pulling \citeyearpar{Pulling2015} and Pulling, Robins, and
Rixon \citeyearpar{Pulling2013} examined long corners in the English Premier
League and demonstrated that deliveries into the goal area and the ``critical
area'' between six and twelve yards from the goal line generate
disproportionately many shot attempts. This finding directy anticipates the
zone system deployed in the present thesis (Section~\ref{sec:data}), which
subdivides the penalty area into depth bands --- Goal Area (GA), Penalty Area
(PA), Penalty Edge (PE), and Penalty Curve (PC) --- and lateral columns. A
recent meta-analysis \citep{JCSS2025} synthesised evidence from 21 studies on
inswinging versus outswinging deliveries and found that outswing corners
generate marginally more final attempts in domestic leagues, attributing the
effect to the ball moving \emph{towards} arriving attackers rather than away
from them. Stafylidis and colleagues \citeyearpar{Stafylidis2022, Stafylidis2023}
reached a similar conclusion but added the qualification that delivery-type
advantage is attacker-dependent: inswing favours late-arriving runners, outswing
favours stationary target players. This interaction between delivery type and
attacker positioning motivates the joint delivery-type / movement-profile
characterisation in SQ2.

The choice of delivery height --- whether the ball arrives flat and driven,
at a medium arc, or as a high-looping floated cross --- has received less
academic attention despite its practical importance. Liu and colleagues
\citeyearpar{Liu2025}, analysing 624 corner-kick goals across the Premier
League, La Liga, and the Danish Superliga alongside 545 World Cup corners,
found significant league-level differences in shot-zone distributions and
inswing-goal proportions, and identified home advantage, inswing delivery,
non-grouped attacking setups, and non-block defensive formations as factors
associated with scoring. Their work also illustrates the methodological
limitation of operating on goal-only samples: by conditioning on goals rather
than all corners, such studies systematically overestimate the effectiveness of
configurations that happened to produce goals while underestimating those that
produced good-but-unconverted chances. The present thesis avoids this by using
the Corner Threat Score (CTS), which integrates delivery quality (xA) with shot
quality (xG) across \emph{all} corners, including those without a subsequent
shot \citep{Lamberts2023}.

\subsubsection{Context, Match Status, and Short Corners}

The notational literature has also catalogued how \emph{context} shapes corner
tactics. Casal, Losada, Maneiro, and Ard\'a \citeyearpar{Casal2017} used CHAID
decision trees on 902 tournament corners to show that match status systematically
alters tactical choices: teams trailing in the final thirty minutes concentrate
attackers closer to the goal and accept greater counter-attack risk, whereas
teams that are level behave more conservatively. De Baranda and
L\'opez-Riquelme \citeyearpar{DeBaranda2012} documented similar effects in the
2006 World Cup. Strafford, Smith, North, and Stone \citeyearpar{Strafford2019},
comparing top-six and bottom-six Premier League sides in 2015/16, reported that
variety of routines did not distinguish the two groups but that execution
consistency did --- an observation that motivates characterising teams by their
\emph{distribution} over routine types rather than by the mere presence of any
given routine. Page and Robins \citeyearpar{PageRobins2012} conducted a
single-club case study in League One and documented substantial within-season
variation in corner tactics, suggesting that clubs actively adjust routines to
opposition, a hypothesis that the per-opponent analysis in SQ5 of this thesis
is designed to test quantitatively.

Short corners --- in which the ball is played to a nearby teammate rather than
crossed directly --- account for roughly 18\,\% of all corner kicks
\citep{Casal2015}, a proportion consistent with the 26\,\% (492 of 1{,}893)
observed in the Dutch Eerste Divisie dataset used in this thesis. Although their
direct goal-conversion rate is lower than for crossed corners, short corners can
serve strategic functions: drawing defenders out of the penalty area before a
subsequent cross, exploiting a numerically weak near-side guard, or recycling
possession when the set-piece structure is unfavourable. Existing academic
literature has characterised short corners primarily at the aggregate level
\citep{Sarmento2025}; the present thesis develops a detailed pass-count and
outcome-based taxonomy (Section~\ref{sec:sq2_short}) that permits more granular
comparison.

\subsubsection{Methodological Limitations of the Notational Tradition}

Two systematic reviews published in 2025 \citep{Sarmento2025, MDPI2025} have
identified persistent methodological limitations in the notational corner-kick
literature. First, \emph{zone inconsistency}: different studies divide the
penalty area into incompatible spatial grids, making cross-study synthesis
unreliable. Second, \emph{outcome operationalisation}: ``corner success'' is
variously defined as a shot, a shot on target, a goal, or the more recent
expected-goals metric, with no consensus standard. Third, \emph{contextual
underdetermination}: most studies report marginal frequencies across categories
without modelling interactions between delivery type, defensive structure, match
status, and player identity. Fourth, and most consequentially for the present
thesis: notational analysis operates on \emph{discrete, human-assigned labels}
and therefore cannot access the continuous spatiotemporal information --- player
velocities, mutual distances, movement correlations --- that tracking data
make available. Manual labels compress rich spatial configurations into single
categorical cells, introducing inter-observer variability and losing the very
information that distinguishes a well-executed zonal block from a poorly
coordinated one.

% ------------------------------------------------------------
\subsection{Defensive and Offensive Set-Piece Systems}
\label{sec:lit_coaching}
% ------------------------------------------------------------

Parallel to the notational academic literature, a body of coaching-science
work focuses on the \emph{systems} teams use to organise themselves at corners.
This literature shares methods with the notational tradition but is oriented
toward design rather than description: rather than asking which variables predict
scoring, it asks how teams should organise and what the trade-offs between
organisational choices are.

\subsubsection{The Zonal--Man--Hybrid Taxonomy}

The foundational distinction in defensive corner organisation is between
\emph{zonal marking} and \emph{man-to-man marking}. In a zonal system, each
defender is assigned a specific area of the penalty box and is responsible for
any ball that enters that area, regardless of which attacker attacks it. In a
man-marking system, each defender is paired with a specific attacker and follows
that attacker's movement until delivery. Pure systems of either type are rare in
elite football: the \emph{hybrid} system, which assigns some defenders to zones
and others to individual attackers, has become the dominant structure at
international and domestic elite level \citep{FIFA2023, MBP2025}.

FIFA's Football Performance Insights report from the 2022 World Cup provides
the most systematic recent taxonomy \citep{FIFA2023}. A defensive setup is
classified as predominantly \emph{zonal} when four or more players are assigned
zonal responsibilities, as predominantly \emph{man} when more than 70\,\% of
the in-box defenders are marking individually, and as \emph{hybrid} otherwise.
Applying this taxonomy to World Cup data, the report found that 26 of 32 teams
brought all ten outfield players back to defend corners, that hybrid setups
dominated, and that strongly zonal teams won a substantially higher proportion
of first contacts on delivered crosses: the United States (five zonal defenders)
cleared 84\,\% of first balls; Germany (six zonal defenders) cleared 75\,\%.

Academic work corroborates these patterns. Pulling and Newton
\citeyearpar{PullingNewton2017} showed that near-post guard systems --- in which
a dedicated defender is stationed at the near post inside or immediately outside
the six-yard box --- materially reduce the probability of a shot from inswinging
corners targeted at that area. Pulling, Robins, and Rixon
\citeyearpar{Pulling2013} documented substantial cross-club variation in how
Premier League sides allocate defenders between zones and man-to-man assignments,
and reported lower concession rates for mixed strategies that included near-post
coverage. Borr\'as and Sainz de Baranda \citeyearpar{Borras2005} and
S\'anchez-Flores \citeyearpar{SanchezFlores2012} reached analogous conclusions
for Spanish football, with the latter also demonstrating that goalkeeper starting
depth relative to the goal line materially changes the spatial coverage of the
zonal block in front of him. The 2025 systematic review \citep{MDPI2025}
synthesises this evidence and concludes that mixed marking strategies with
near-post coverage consistently yield lower concession rates than either pure
alternative, while acknowledging that the inconsistency of zone definitions
across studies makes precise quantification difficult.

The coaching literature is explicit about the trade-offs. Zonal marking
concentrates defenders between the ball and goal, simplifies communication, and
requires defenders to \emph{attack the ball} rather than track a moving player;
but it is vulnerable to deliberate blocking runs that occupy a zone without
attacking the ball, and to deliveries aimed precisely between two zonal
defenders. Man-marking eliminates the problem of uncontested runners and allows
defenders' physical attributes to be matched against specific threats, but is
\emph{reactive} by construction: a man-marker who watches his attacker rather
than the ball is perpetually a step behind on delivery \citep{TacticalAnalyst2016}.
Individual defensive errors propagate into goals because there is no positional
safety net. The hybrid compromise reduces the severity of both failure modes but
introduces new ones: confusion about responsibility when attackers move between
zones and markers, and difficulty adjusting to short corners that require the
zonal block to reorganise quickly. This three-way taxonomy is directly reflected
in the defensive strategy classification developed for SQ1 of this thesis
(Section~\ref{sec:sq1_team}), which distinguishes setups by marking intensity
and the presence or absence of short and counter defenders.

\subsubsection{Specific Defensive Roles}

Beyond the aggregate system label, the coaching literature identifies specific
\emph{positional roles} that individual defenders can occupy during a corner.
These include: the near-post guard (stationed inside the six-yard box at the
near post); the far-post guard (analogous position at the far post); the
short-corner guard (stationed close to the ball to press a short corner option);
the back-space defender (positioned at the edge of the penalty area to prevent
cut-back shots and cover the counter-attack); and the man-marker (paired with a
specific attacker). MBP School \citeyearpar{MBP2025} describes a typical hybrid
setup as employing two or three zonal defenders between the six-yard box and the
penalty spot, one near-post guard, two to three man-markers on the most dangerous
aerial targets, and one or two edge defenders. This granular role vocabulary is
directly adopted in the seven-role taxonomy of Bauer, Anzer, and Smith
\citeyearpar{Bauer2022} and is reflected in the four-role classification
(MAN, ZONAL, SHORT, COUNTER) used in SQ1 of this thesis, which collapses the
full taxonomy to achieve classification accuracy on the available labelled
dataset.

\subsubsection{Attacking Routines and Tactical Design}

On the attacking side, the coaching and applied practitioner literature has
moved considerably ahead of the academic literature in recent years. Early
academic work catalogued a binary distinction between \emph{static} offensive
organisation (players start from fixed positions) and \emph{dynamic} organisation
(players begin grouped and execute rehearsed runs) \citep{Maneiro2014}. Maneiro's
analysis found that dynamic organisation is associated with higher shot-conversion
rates, a finding replicated in women's elite football by Beare and Stone
\citeyearpar{BeareStone2019}. This static/dynamic dimension is directly reflected
in the movement-profile classification of SQ2 (Section~\ref{sec:sq2_runs}).

Practitioner writing has developed a much richer vocabulary than the academic
literature captures. Set-piece specialists --- most notably Gianni Vio, whose
work was first operationalised at Brentford and later used in the Italian
national team's UEFA Euro 2020 campaign, and Nicolas Jover, whose approach
transformed Arsenal's dead-ball effectiveness --- describe attacking corner
routines in terms of \emph{themes}: coordinated multi-player configurations
with specific strategic intentions. Common themes include \emph{stacking}
(clustering three or more attackers at the near post to force defensive
overcommitment, then breaking in multiple directions), \emph{cross-stacking}
(a near-post stack positioned on the far side, inviting cross-goal runs),
\emph{box} setups (four attackers forming a rectangle that disperses in four
directions on delivery), \emph{train} setups (attackers in single file, splitting
at delivery), and \emph{targets-and-screens} routines (in which one attacker's
primary role is to physically obstruct a specific defender's movement so that a
second attacker can receive the ball unmarked). Each of these themes is defined
not by any single player's action but by the \emph{coordinated pattern} of
multiple players' starting positions and runs --- precisely the kind of
multi-player spatiotemporal pattern that tracking data can capture but
categorical notational codes cannot. The goal of SQ2 in this thesis is to
operationalise this practitioner vocabulary as a data-derived taxonomy of
attacking movement profiles, grounded in manual observation of 86 corners and
automated via a rule-based running-line classifier applied to the full
1{,}893-corner dataset.

% ------------------------------------------------------------
\subsection{Tracking Data and Machine Learning Applied to Corner Kicks}
\label{sec:lit_tracking}
% ------------------------------------------------------------

The application of tracking data to football tactics has matured rapidly since
Bialkowski and colleagues \citeyearpar{Bialkowski2014, Bialkowski2016}
demonstrated that spatiotemporal player trajectories could be used to
automatically recover team formations and individual positional roles. Rein and
Memmert \citeyearpar{ReinMemmert2016} established the theoretical framework for
big-data approaches to football performance, arguing that the central challenge
is not data availability but the development of theoretically grounded models
that translate high-frequency positional observations into tactically meaningful
constructs. Goes and colleagues \citeyearpar{Goes2019, Goes2021} subsequently
demonstrated that features derived from tracking data --- relative velocities,
distance to nearest opponents, time-to-goal measures, and convex hull area
--- can be used to model pass effectiveness, team performance, and space
creation, establishing a feature-engineering vocabulary that has since been
widely adopted. Recent systematic reviews \citep{RicoGonzalez2023,
Clemente2025, OttingKarlis2023} document a proliferation of methods applied to
the open-play problem: hidden Markov models for tactical-state identification,
graph neural networks for multi-agent interaction modelling, Delaunay
triangulations for formation analysis, and unsupervised clustering methods for
discovering latent tactical structure.

Within this broader field, the set-piece-specific literature is still emerging
but has produced four landmark contributions that are in direct methodological
dialogue with this thesis.

\subsubsection{Power et al.\ (2018): Mythbusting Set Pieces}

Power, Hobbs, Ruiz, Wei, and Lucey \citeyearpar{Power2018} introduced the first
large-scale, tracking-data-based analysis of set pieces in \emph{Mythbusting
Set-Pieces in Soccer}, presented at the MIT Sloan Sports Analytics Conference.
Working with Premier League data from the 2016/17 season, they demonstrated that
set-piece possessions generate goals at 1.8\,\% per possession versus 1.1\,\% for
open-play possessions, providing the quantitative grounding for treating set
pieces as a market inefficiency. Their paper introduced image-based
representations of set-piece configurations: the positions of all players on
the pitch were encoded as spatial heatmaps, and image-similarity measures were
used to retrieve corners from a large database that matched a given target
situation. Critically, they framed the defensive-side analysis as an explicit
search for exploitable inefficiencies --- gaps between what a defence is designed
to prevent and what attacking routines actually challenge it with --- setting an
agenda that subsequent work, including this thesis, has continued to pursue.

\subsubsection{Shaw and Gopaladesikan (2020): The Corner-Kick Playbook}

Shaw and Gopaladesikan's \emph{Routine Inspection: A Playbook for Corner Kicks}
\citeyearpar{ShawGopaladesikan2020} is the most methodologically influential
prior work for this thesis, and the winner of the research paper competition at
the 2021 MIT Sloan Sports Analytics Conference. It is the first study to address
both the offensive and defensive characterisation of corners from tracking data
within a single framework, and it establishes the \emph{player-level} approach
to set-piece analysis that is adopted and extended here.

On the offensive side, they defined an attacking \emph{run} as a player's
starting location two seconds before the corner is taken and their \emph{target
location} either one second after first ball contact or two seconds after
delivery, whichever comes first. Critically, the ball trajectory is excluded from
this definition: the intended target of a cross frequently differs from where the
ball arrives, making the ball endpoint an unreliable proxy for attacking intent.
They used Gaussian mixture modelling to cluster starting locations into six zones
and target locations into seven zones, producing 42 theoretically possible
player runs; non-negative matrix factorisation (NMF) then compressed
co-occurring runs into 30 ``routine'' basis vectors. An entire team's season of
corners could thereby be summarised as a distribution over routines --- a
statistical playbook for opposition scouting. On the defensive side, they trained
a supervised classifier on manually labelled corners to distinguish zonal from
man-marking defenders at the player level, achieving approximately 81\,\% precision
and recall for zonal classification. Their player-level aggregation approach ---
characterising a team's defensive setup as the ratio and spatial distribution of
individually classified roles --- is directly adopted in SQ1 of this thesis.

\subsubsection{Bauer, Anzer, and Smith (2022): Fine-Grained Role Classification}

Bauer, Anzer, and Smith \citeyearpar{Bauer2022} extended the defensive
classification problem to a richer taxonomy. Rather than a binary zonal/man
distinction, they identified seven defensive roles: player-marking, zonal-marking,
positioned for counterattack, back-space defender, short-corner guard, near-post
defender, and far-post defender. Working with Bundesliga tracking data, they
hand-labelled the role of every defensive player in 213 corners across 33 matches
and trained an ensemble classifier with data augmentation, achieving a weighted
accuracy of 89.3\,\% against inter-labeller ground truth. Crucially, the model
correctly identified the specific attacker being man-marked 80.8\,\% of the time,
demonstrating that tracking data encode sufficient information to solve the
\emph{assignment} problem --- not just the role-classification problem. They also
demonstrated that a rule-based baseline could be used as weak supervision to
bootstrap learning on unlabelled data, following the Snorkel paradigm
\citep{Ratner2017}.

Three aspects of their methodology are directly adopted in SQ1 of this thesis.
First, the two-step pipeline: step one classifies each defender's role (MAN,
ZONAL, SHORT, or COUNTER); step two applies the Hungarian assignment algorithm
to assign specific attacking targets to the identified man-markers as a
minimum-cost bipartite matching problem. Second, the feature engineering approach:
pair-wise distance and velocity features computed over multiple pre-delivery
time windows (2, 3, and 5 seconds), motion correlation between defender and
attacker trajectories, mutual nearest-neighbour indicators, and distance to
ball. Third, the leave-one-match-out cross-validation scheme, which prevents
information leakage between training and test sets at the match level.

\subsubsection{Wang et al.\ (2024): TacticAI}

Wang and colleagues \citeyearpar{Wang2024} published \emph{TacticAI} in
\emph{Nature Communications}, representing the most technically ambitious
application of machine learning to corner-kick analysis to date. TacticAI was
trained on 9{,}693 corners from three Premier League seasons in collaboration
with Liverpool FC. It represents each corner as a graph in which nodes are
players (with positions and velocities as node features) and edges encode spatial
relationships, and uses geometric graph neural networks --- invariant to
pitch-symmetry transformations --- to perform three tasks: predicting the
receiver of the corner delivery, predicting whether a shot attempt occurs, and
\emph{generating} suggested attacker repositioning to maximise the probability
of a shot. In a blinded evaluation with Liverpool's coaching and analysis staff,
TacticAI's suggested player setups were judged preferable to the original
configuration in approximately 90\,\% of cases. This result demonstrates that
deep learning on tracking-derived graph representations can produce tactically
credible and practically useful outputs.

TacticAI is a landmark contribution and its architectural choices --- geometric
deep learning, graph-based player representations, symmetry invariance --- are
likely to define the next generation of set-piece analytics tools. However, its
scope is confined to the \emph{prescriptive} question on the offensive side:
given an observed defensive setup, where should attackers stand? It does not
produce an interpretable, separable characterisation of the defensive system
itself, it does not classify attacking routines into a taxonomy that coaching
staff can design from, and it does not address the conditional effectiveness
question --- which attacking approach is most effective against which defensive
type --- that is the central objective of this thesis. The present work is
therefore complementary to TacticAI rather than competing with it: it produces
the interpretable typologies of both offensive and defensive behaviour that make
the conditional effectiveness question answerable in the language that coaches
use.

% ------------------------------------------------------------
\subsection{Individual Player Attributes and Aerial Performance}
\label{sec:lit_players}
% ------------------------------------------------------------

A fourth, more dispersed body of literature concerns the contribution of
individual player attributes to corner outcomes. The aerial duel is the primary
mechanism by which corner deliveries are converted into shots or cleared. An
aerial duel is operationally defined by providers such as Wyscout and Opta as a
contest in which two or more opposing players jump simultaneously for a ball in
the air; the duel is won by the player who touches the ball first. Derived
statistics --- aerial duels won per 90 minutes, aerial duel win rate --- are
standard features in player scouting platforms and are widely used as proxies
for aerial ability in applied work.

The physical and technical determinants of aerial performance are multi-dimensional.
Forte and colleagues \citeyearpar{Forte2019}, in an observational analysis of
920 headers from the 2017/18 Premier League, showed that set-piece situations
account for a substantial proportion of all heading actions, that most headers
result in possession being lost rather than retained, and that pre-jump body
orientation, contact surface, and positioning relative to the ball explain a
meaningful share of heading outcome variance. Height, jump height, and timing
are the three attributes most consistently cited by coaches in assessing aerial
threat; however, they are only loosely correlated with raw aerial duel win rate
because win rate conflates a player's own ability with the quality and height of
the specific opponents they have been tested against.

Kim, Lee, and Park \citeyearpar{Kim2024}, addressing precisely this confound,
applied a height-adjusted Elo rating model to all aerial duels across three
seasons of the Korean K League 1. They found that player rankings by adjusted
Elo differed substantially from rankings by raw win rate: the player ranked first
by Elo placed third by raw rate, and several players in the raw top-twenty fell
out of the adjusted top-fifty entirely. This result has direct implications for
SQ4 of this thesis: any attempt to link corner outcomes to individual aerial
attributes must use opponent-adjusted or context-normalised metrics rather than
raw duel win rates, which systematically favour players who happened to face
weaker aerial opponents rather than those with genuinely superior aerial
technique. In the absence of an Elo-based estimate for all players in the Eerste
Divisie dataset, this thesis uses aerial duel win rate as the primary attribute
while acknowledging its limitations, and cross-references with player height
where metadata are available.

The existing corner-kick literature has engaged only lightly with individual
attributes. Casal and colleagues \citeyearpar{Casal2015, Casal2017}, Pulling and
Newton \citeyearpar{PullingNewton2017}, and most notational studies treat players
as interchangeable within their tactical roles. Shaw and Gopaladesikan
\citeyearpar{ShawGopaladesikan2020} and Bauer, Anzer, and Smith
\citeyearpar{Bauer2022} do not incorporate individual attributes in their
classification models. Wang et al.\ \citeyearpar{Wang2024} use position and
velocity as node features in TacticAI but do not attempt to model player-specific
aerial skill. The player-quality dimension has therefore been characterised and
theoretically motivated but not integrated with the structural analysis of corner
tactics, leaving the interaction between player ability and tactical setup ---
the question that distinguishes a decision-support tool from a population-level
statistical summary --- unaddressed in the existing literature.

% ------------------------------------------------------------
\subsection{Synthesis: The Research Gap}
\label{sec:lit_gap}
% ------------------------------------------------------------

The four bodies of literature reviewed above each make important contributions,
but each has a boundary that prevents it from answering the question posed in
Section~\ref{sec:research_questions}. The notational tradition provides a rich
descriptive vocabulary and a well-established set of performance indicators, but
operates on discrete categorical variables and cannot access the continuous
spatiotemporal information that distinguishes tactical intent from spatial
outcome. The coaching-science literature provides the conceptual vocabulary ---
zonal, man, hybrid, stacking, blocking, targeting --- that coaches use in
practice, but does not derive these concepts automatically from data and does not
evaluate them conditionally on both sides of the corner simultaneously. The
tracking-data literature has demonstrated that individual defensive roles can be
classified from positional data with high accuracy \citep{Bauer2022}, that
attacking routines can be represented as distributions over player-movement
patterns \citep{ShawGopaladesikan2020}, and that deep learning can generate
tactically credible positional suggestions \citep{Wang2024} --- but no study has
combined a data-derived characterisation of both sides into a single conditional
effectiveness framework. The individual-attributes literature has shown that
raw aerial duel win rates are insufficient proxies for aerial ability
\citep{Kim2024} and that heading outcome is multidimensionally determined
\citep{Forte2019}, but has not been integrated with the structural analysis of
set-piece tactics.

The gap is therefore not a lack of data, methods, or descriptive knowledge.
It is the absence of an integrated, data-driven framework that simultaneously
characterises the attacking routine, the defensive structure, and the individual
player attributes involved in a given corner kick, and then uses these
characterisations jointly to answer the question that motivates all set-piece
coaching: \emph{which of our attacking approaches is most likely to succeed
against this opponent's specific defensive configuration, given our players'
qualities?}

This thesis builds that framework. Using a dataset of 1{,}893 corner kicks from
the 2025/26 Dutch Eerste Divisie season, supplemented by 86 manually labelled
corners that serve as ground truth, the following chapters develop a five-layer
analytical pipeline --- defensive role classification (SQ1), offensive pattern
classification (SQ2), conditional effectiveness analysis (SQ3), player-quality
modelling (SQ4), and a practical preparation tool (SQ5) --- that moves from raw
tracking and event data to opponent-specific tactical recommendations. The
methodological design of this pipeline is described in full in
Chapter~\ref{sec:methods}.